\section{Introdução}

A inteligência artificial (IA) e o aprendizado de máquina (\textit{machine learning} --- ML) constituem o substrato cognitivo inalienável dos gêmeos digitais odontológicos contemporâneos. Enquanto os capítulos precedentes delimitaram os alicerces da modelagem geométrica e biomecânica (\autoref{ch:fundamentos}), além do panorama evolutivo da odontologia digital (\autoref{ch:panorama}), o presente escopo dedica-se a dissecar os algoritmos e as arquiteturas subjacentes que dotam os gêmeos digitais de suas prerrogativas de aprendizado contínuo, inferência preditiva e autopoiese computacional.

Desprovido de IA, um pretenso gêmeo digital odontológico não passaria de uma virtualização tridimensional inerte --- uma miragem topológica, rigorosa em métricas, porém vazia de agentividade. É a intrincada malha de IA que outorga ao gêmeo digital a autonomia para segmentar milimetricamente estruturas anatômicas em tomografias computadorizadas de feixe cônico (CBCT), prever trajetórias ortodônticas sob pressões vetoriais complexas, parametrizar implantes mediante simulações de elementos finitos e identificar estigmas patológicos sub-clínicos frequentemente elididos pelo crivo humano.

O interstício temporal entre 2024 e 2026 testemunhou uma maturação hiperbólica da IA odontológica. Este salto não foi fortuito; derivou de uma tríade convergente: a democratização de imensos *datasets* clínicos heterogêneos, o refinamento das topologias de aprendizado profundo (\textit{deep learning}) --- em particular, os Transformers --- e o advento de hardware especializado e acessível. Conforme elucidado por Katsoulakis et al. (2024) em ampla revisão sobre gêmeos digitais na saúde \cite{katsoulakis2024digital}, a orquestração multinível promovida pela IA desloca o paradigma da representação estática para simulações biossistêmicas interativas. Na seara odontológica, Duggal et al. (2026) ressaltam o potencial incipiente, mas disruptivo, da tecnologia para simulação e monitoramento hiperpersonalizado \cite{duggal2026exploring}. Da mesma forma, Roy et al. (2025) advogam que essa amálgama transfigura irreparavelmente o diagnóstico e a intervenção \cite{roy2025editorial}.

\destaque{O amadurecimento dos gêmeos digitais na odontologia demarca a transição incontestável de um modelo reativo-reparador para uma práxis preventiva, preditiva e de altíssima precisão empírica.}

Este capítulo organiza-se de maneira a esgotar as facetas mais prementes do tema. Inicia-se delineando a epistemologia e os alicerces teóricos da IA na disciplina. Subsequente a isso, as seções avançam sobre as aplicações fáticas já em estado da arte, a necessidade inegociável de explicabilidade algorítmica (\textit{Explainable AI} --- XAI), o impacto pervasivo dos Grandes Modelos de Linguagem (LLMs), a descoberta de biomarcadores suportados por \textit{deep learning} e, não menos importante, a dissecação crítica dos gargalos metodológicos, regulatórios e socioeconômicos que ainda obnubilam sua ampla capilaridade no sistema de saúde.
